% !TEX root = ../main.tex

\chapter{数值实验}\label{ch:numerical}
  本章用两个线性方程和两个半线性方程检验前文方法.
  线性方程包括 Black--Scholes 方程和 Heat 方程, 用来验证线性最小二乘求解器;
  半线性方程包括 HJB-LQ 方程和 Allen--Cahn 方程, 用来考察 Levenberg--Marquardt 方法在高维 BSDE 离散系统中的表现.

  \section{实验设置}\label{sec:numerical_setup}
    对每个方程, 先生成样本路径
    $\{X_k^{(i)},\Delta W_k^{(i)}\}_{k=0}^{N-1}$,
    再用随机特征函数近似各时间层上的 $Z_k$.
    线性方程中, 终端值 $Y_N$ 关于 $y_0$ 和输出层参数 $\alpha$ 是线性的,
    因而直接求解带 ridge 正则的最小二乘问题.
    半线性方程中, 终端残差关于参数不再线性, 因而采用 LM 方法求解非线性最小二乘问题.

    本章统一报告训练样本上的终端残差 RMSE 和独立测试样本上的终端残差 RMSE\@.
    前者反映固定样本路径上的拟合程度, 后者用于判断随机特征参数是否发生了路径级过拟合.
    除特别说明外, 标准配置见表~\ref{tab:numerical_config}.

    \begin{table}
      \centering
      \caption{主要实验配置}
      \label{tab:numerical_config}
      \begin{tabular}{lcccccccc}
        \toprule
        方程 & 类型 & $d$ & $T$ & $N$ & $S$ & $H$ & $\lambda_0$ & 求解器 \\
        \midrule
        Black--Scholes & 线性 & 100 & 1.0 & 100 & 16384 & 50 & 0.010 & 线性最小二乘 \\
        Heat & 线性 & 100 & 1.0 & 100 & 16384 & 50 & 0.010 & 线性最小二乘 \\
        HJB-LQ & 半线性 & 100 & 1.0 & 20 & 16384 & 50 & 0.010 & LM \\
        Allen--Cahn & 半线性 & 100 & 0.3 & 20 & 16384 & 50 & 0.010 & LM \\
        \bottomrule
      \end{tabular}
    \end{table}

  \section{线性方程}\label{sec:numerical_linear}
    本节考虑生成元关于 $(y,z)$ 具有线性结构的方程.
    对这类问题, 离散 BSDE 的终端值可以写成
    \begin{equation}
      Y_N^{(i)}
      = A_{i,0} y_0
      + \sum_{j=1}^{dH} A_{i,j} \theta_j
      + b_i,
      \qquad
      \theta=\operatorname{vec}(\alpha).
      \label{eq:numerical_linear_terminal_form}
    \end{equation}
    因此终端匹配条件 $Y_N^{(i)}\approx g(X_N^{(i)})$ 给出一个线性最小二乘问题,
    不需要外层非线性迭代.

    Black--Scholes 算例采用多资产算术平均看涨期权.
    状态过程满足
    \begin{equation}
      X_{k+1}^{m}
      =
      X_k^{m}\exp\!\left((r-\frac{1}{2}\sigma^2)\Delta t
      + \sigma \Delta W_k^{m}\right),
      \qquad m=1,\ldots,d,
      \label{eq:numerical_bsm_forward}
    \end{equation}
    生成元和终端条件为
    \begin{equation}
      f(t,x,y,z)=-r y,\qquad
      g(x)=\max\!\left(\frac{1}{d}\sum_{m=1}^{d}x_m-K,0\right).
      \label{eq:numerical_bsm_terminal}
    \end{equation}
    由于该 payoff 没有简单闭式价格, 本文采用控制变量 Monte Carlo 参考值, 见附录~\ref{ch:reference_values}.

    Heat 算例采用
    \begin{equation}
      X_t=x_0+W_t,\qquad
      f(t,x,y,z)=0,\qquad
      g(x)=\exp\!\left(-\frac{\lVert x\rVert^2}{a}\right).
      \label{eq:numerical_heat}
    \end{equation}
    由 Feynman--Kac 公式,
    \begin{equation}
      u(t,x)
      =
      \left(\frac{a}{a+2(T-t)}\right)^{d/2}
      \exp\!\left(
        -\frac{\lVert x\rVert^2}{a+2(T-t)}
      \right).
      \label{eq:numerical_heat_solution}
    \end{equation}
    因而在 $x_0=0$, $a=d$ 时,
    \begin{equation}
      Y_0=u(0,0)=\left(\frac{d}{d+2T}\right)^{d/2}.
      \label{eq:numerical_heat_reference}
    \end{equation}

    \begin{table}
      \centering
      \caption{线性方程主要结果}
      \label{tab:linear_main_results}
      \begin{tabular}{lccccc}
        \toprule
        方程 & $y_0$ & 参考值 & 绝对误差 & 训练 RMSE & 测试 RMSE \\
        \midrule
        Black--Scholes & 0.048701 & 0.048809 & $1.08\times10^{-4}$ & 0.0112 & 0.0158 \\
        Heat & 0.371605 & 0.371528 & $7.71\times10^{-5}$ & 0.0356 & 0.0529 \\
        \bottomrule
      \end{tabular}
    \end{table}

    由表~\ref{tab:linear_main_results} 可见, 两个线性算例的初值都与独立参考值较为接近.
    Black--Scholes 的误差约为 $10^{-4}$, Heat 方程的误差小于 $10^{-4}$.

  \section{半线性方程}\label{sec:numerical_nonlinear}
    HJB-LQ 方程的非线性来自梯度项, Allen--Cahn 方程的非线性来自反应项.
    这两类非线性分别检验方法对 $z$ 型非线性和 $y$ 型非线性的处理能力.

    HJB-LQ 方程采用
    \begin{equation}
      X_t=\sqrt{2}W_t,\qquad
      f(t,x,y,z)=-\frac{1}{2}\lVert z\rVert^2,
      \qquad
      g(x)=\log\!\left(\frac{1+\lVert x\rVert^2}{2}\right).
      \label{eq:numerical_hjblq}
    \end{equation}
    对应 BSDE 为
    \begin{equation}
      \dif Y_t=\frac{1}{2}\lVert Z_t\rVert^2\dif t+Z_t^{\operatorname{T}}\dif W_t.
      \label{eq:numerical_hjblq_bsde}
    \end{equation}
    令 $U_t=\exp(-Y_t)$, 可得 $U_t$ 为鞅.
    因此
    \begin{equation}
      Y_0
      =
      -\log \mathbb{E}
      \left[
        \frac{2}{1+\lVert X_T\rVert^2}
      \right]
      =
      -\log \mathbb{E}
      \left[
        \frac{2}{1+2T\chi_d^2}
      \right].
      \label{eq:numerical_hjblq_reference}
    \end{equation}
    当 $d=100$, $T=1$ 时, 一维积分参考值为 $Y_0=4.5901617246$.

    Allen--Cahn 方程采用
    \begin{equation}
      X_t=\sqrt{2}W_t,\qquad
      f(t,x,y,z)=y-y^3,
      \qquad
      g(x)=\frac{1}{2+0.4\lVert x\rVert^2}.
      \label{eq:numerical_allencahn}
    \end{equation}
    该方程没有简单闭式解.
    本文采用独立的径向有限差分方法构造参考值.
    对 $d=100$, $T=0.3$ 的配置, 参考值为 $Y_0=0.05278464$.

    \begin{table}
      \centering
      \caption{半线性方程主要结果}
      \label{tab:nonlinear_main_results}
      \begin{tabular}{lccccc}
        \toprule
        方程 & $y_0$ & 参考值 & 绝对误差 & 训练 RMSE & 测试 RMSE \\
        \midrule
        HJB-LQ & 4.590000 & 4.590162 & $1.62\times10^{-4}$ & 0.0990 & 0.1443 \\
        Allen--Cahn & 0.053131 & 0.052785 & $3.46\times10^{-4}$ & 0.0035 & 0.0051 \\
        \bottomrule
      \end{tabular}
    \end{table}

    表~\ref{tab:nonlinear_main_results} 表明, LM 求解后的两个半线性算例都能给出较准确的初值.
    HJB-LQ 的初值误差约为 $1.62\times10^{-4}$,
    Allen--Cahn 的初值误差约为 $3.46\times10^{-4}$.

    \begin{table}
      \centering
      \caption{半线性方程的 LM 迭代结果}
      \label{tab:nonlinear_iteration_results}
      \begin{tabular}{lccccc}
        \toprule
        方程 & 迭代步数 & 接受步数 & 最终 $\lambda$ & 最终步长范数 & 时间 / ms \\
        \midrule
        HJB-LQ & 4 & 3 & $1.00\times10^{4}$ & $2.82\times10^{-7}$ & 3435.640 \\
        Allen--Cahn & 4 & 3 & $1.00\times10^{6}$ & $5.59\times10^{-10}$ & 3145.050 \\
        \bottomrule
      \end{tabular}
    \end{table}

    图~\ref{fig:nonlinear_residual_history} 展示半线性求解器的残差下降过程.
    横轴为 LM 接受步编号, 纵轴为训练终端残差 RMSE\@.

    \begin{figure}[ht]
      \centering
      \includegraphics[width=0.82\textwidth]{nonlinear-lm-residual}
      \caption{半线性方程的 LM 迭代残差下降}
      \label{fig:nonlinear_residual_history}
    \end{figure}

  \section{样本数的影响}\label{sec:numerical_sample_size}
    样本数直接影响最小二乘系统对真实分布的代表性.
    表~\ref{tab:sample_size_results} 固定 $d=100$, $H=50$, 只改变训练样本数 $S$.

    \begin{table}
      \centering
      \caption{样本数变化结果}
      \label{tab:sample_size_results}
      \begin{tabular}{lccccc}
        \toprule
        方程 & $S$ & $y_0$ & 绝对误差 & 训练 RMSE & 测试 RMSE \\
        \midrule
        Heat & 4096 & 0.373519 & $1.99\times10^{-3}$ & 0.0002 & 0.1487 \\
        Heat & 8192 & 0.372477 & $9.49\times10^{-4}$ & 0.0268 & 0.0700 \\
        Heat & 16384 & 0.371605 & $7.71\times10^{-5}$ & 0.0356 & 0.0529 \\
        Black--Scholes & 8192 & 0.040108 & $8.70\times10^{-3}$ & 0.6489 & 1.5086 \\
        Black--Scholes & 16384 & 0.048701 & $1.08\times10^{-4}$ & 0.0112 & 0.0158 \\
        HJB-LQ & 4096 & 4.423910 & $1.66\times10^{-1}$ & $1.36\times10^{-6}$ & 0.4595 \\
        HJB-LQ & 8192 & 4.578700 & $1.15\times10^{-2}$ & 0.0762 & 0.1963 \\
        HJB-LQ & 16384 & 4.590000 & $1.62\times10^{-4}$ & 0.0990 & 0.1443 \\
        Allen--Cahn & 8192 & 0.053268 & $4.84\times10^{-4}$ & 0.0028 & 0.0069 \\
        Allen--Cahn & 16384 & 0.053131 & $3.46\times10^{-4}$ & 0.0035 & 0.0051 \\
        \bottomrule
      \end{tabular}
    \end{table}

    小样本时训练残差不一定可靠.
    例如 HJB-LQ 在 $S=4096$ 时训练 RMSE 几乎为零, 但初值误差和测试 RMSE 都明显偏大.
    因此后续判断以独立测试残差和初值误差为主.
    图~\ref{fig:sample_size_generalization} 给出了相同数据的图形化比较.

    \begin{figure}[ht]
      \centering
      \includegraphics[width=0.88\textwidth]{sample-size-generalization}
      \caption{样本数对测试残差和初值误差的影响}
      \label{fig:sample_size_generalization}
    \end{figure}

  \section{随机特征宽度的影响}\label{sec:numerical_hidden_dim}
    随机特征宽度 $H$ 控制每个时间层上可训练输出参数的数量.
    表~\ref{tab:hidden_dim_results} 固定 $d=100$, $S=16384$, 只改变 $H$.

    \begin{table}
      \centering
      \caption{随机特征宽度变化结果}
      \label{tab:hidden_dim_results}
      \begin{tabular}{lccccc}
        \toprule
        方程 & $H$ & $y_0$ & 绝对误差 & 训练 RMSE & 测试 RMSE \\
        \midrule
        Heat & 8 & 0.372033 & $5.05\times10^{-4}$ & 0.0491 & 0.0521 \\
        Heat & 20 & 0.371917 & $3.89\times10^{-4}$ & 0.0454 & 0.0520 \\
        Heat & 50 & 0.371605 & $7.71\times10^{-5}$ & 0.0356 & 0.0529 \\
        Black--Scholes & 20 & 0.048932 & $1.23\times10^{-4}$ & 0.0140 & 0.0159 \\
        Black--Scholes & 50 & 0.048701 & $1.08\times10^{-4}$ & 0.0112 & 0.0158 \\
        HJB-LQ & 8 & 4.597250 & $7.09\times10^{-3}$ & 0.1353 & 0.1429 \\
        HJB-LQ & 12 & 4.596820 & $6.66\times10^{-3}$ & 0.1321 & 0.1433 \\
        HJB-LQ & 20 & 4.594760 & $4.60\times10^{-3}$ & 0.1251 & 0.1430 \\
        HJB-LQ & 50 & 4.590000 & $1.62\times10^{-4}$ & 0.0990 & 0.1443 \\
        Allen--Cahn & 20 & 0.053057 & $2.72\times10^{-4}$ & 0.0045 & 0.0052 \\
        Allen--Cahn & 50 & 0.053131 & $3.46\times10^{-4}$ & 0.0035 & 0.0051 \\
        \bottomrule
      \end{tabular}
    \end{table}

    对 HJB-LQ 方程, 增大 $H$ 明显改善初值误差和训练残差;
    对 Heat, Black--Scholes 和 Allen--Cahn, 当前测试误差对 $H$ 的变化较平缓.
    图~\ref{fig:hidden_dim_generalization} 展示了测试残差和初值误差随 $H$ 的变化.

    \begin{figure}[ht]
      \centering
      \includegraphics[width=0.88\textwidth]{hidden-dim-generalization}
      \caption{随机特征宽度对测试残差和初值误差的影响}
      \label{fig:hidden_dim_generalization}
    \end{figure}

  \section{维度与运行时间}\label{sec:numerical_dimension_time}
    为观察维度变化对运行时间的影响, 本节固定 $H=50$, $S=16384$,
    对 $d=20,50,100$ 分别运行三个随机种子并取平均.
    表~\ref{tab:dimension_average_results} 给出平均结果.

    \begin{table}
      \centering
      \caption{维度变化的平均结果}
      \label{tab:dimension_average_results}
      \begin{tabular}{lcccc}
        \toprule
        方程 & $d$ & $y_0$ 均值 & 测试 RMSE 均值 & 时间 / ms \\
        \midrule
        Heat & 20 & 0.385587 & 0.0500 & $299.2\pm1.2$ \\
        Heat & 50 & 0.375873 & 0.0528 & $434.1\pm1.4$ \\
        Heat & 100 & 0.371924 & 0.0520 & $697.3\pm2.6$ \\
        Black--Scholes & 20 & 0.051706 & 0.0356 & $301.1\pm1.9$ \\
        Black--Scholes & 50 & 0.049297 & 0.0390 & $440.9\pm2.1$ \\
        Black--Scholes & 100 & 0.048774 & 0.0124 & $704.6\pm0.7$ \\
        HJB-LQ & 20 & 2.934457 & 0.1535 & $447.9\pm70.8$ \\
        HJB-LQ & 50 & 3.888443 & 0.1499 & $1039.8\pm181.7$ \\
        HJB-LQ & 100 & 4.591863 & 0.1436 & $3324.4\pm265.6$ \\
        Allen--Cahn & 20 & 0.206768 & 0.0154 & $442.8\pm58.4$ \\
        Allen--Cahn & 50 & 0.099445 & 0.0088 & $858.8\pm113.5$ \\
        Allen--Cahn & 100 & 0.053021 & 0.0051 & $2764.5\pm652.0$ \\
        \bottomrule
      \end{tabular}
    \end{table}

    用幂律模型 $t(d)\approx C d^p$ 对平均时间做最小二乘拟合,
    得到 Heat 的经验阶数 $p=0.52$, Black--Scholes 的经验阶数 $p=0.52$,
    HJB-LQ 的经验阶数 $p=1.23$, Allen--Cahn 的经验阶数 $p=1.12$.
    线性方程在这组维度上增长较慢, 半线性方程因为需要重复构造残差和灵敏度矩阵, 时间增长更明显.
    图~\ref{fig:dimension_time_scaling} 在对数坐标下展示了这一趋势.

    \begin{figure}[ht]
      \centering
      \includegraphics[width=0.82\textwidth]{dimension-time-scaling}
      \caption{运行时间随维度变化的经验阶数}
      \label{fig:dimension_time_scaling}
    \end{figure}

  \section{随机种子稳定性}\label{sec:numerical_seed}
    随机特征内层参数和样本路径都由随机数生成, 因此需要检查单次结果是否依赖某个特殊种子.
    表~\ref{tab:seed_stability} 使用标准配置下的三个种子,
    报告 $y_0$ 均值, 样本标准差和平均测试残差.

    \begin{table}
      \centering
      \caption{随机种子稳定性结果}
      \label{tab:seed_stability}
      \begin{tabular}{lcccc}
        \toprule
        方程 & $y_0$ 均值 & $y_0$ 标准差 & 均值绝对误差 & 测试 RMSE 均值 \\
        \midrule
        Heat & 0.371924 & $2.81\times10^{-4}$ & $3.96\times10^{-4}$ & 0.0520 \\
        Black--Scholes & 0.048774 & $6.30\times10^{-5}$ & $3.56\times10^{-5}$ & 0.0124 \\
        HJB-LQ & 4.591863 & $1.62\times10^{-3}$ & $1.70\times10^{-3}$ & 0.1436 \\
        Allen--Cahn & 0.053021 & $9.54\times10^{-5}$ & $2.37\times10^{-4}$ & 0.0051 \\
        \bottomrule
      \end{tabular}
    \end{table}

    四个方程的 $y_0$ 波动均小于主要数值量级.
    图~\ref{fig:seed_stability} 给出每个随机种子下的初值绝对误差.

    \begin{figure}[ht]
      \centering
      \includegraphics[width=0.78\textwidth]{seed-stability}
      \caption{不同随机种子下的初值误差}
      \label{fig:seed_stability}
    \end{figure}

  \section{与 DeepBSDE 方法的对比}\label{sec:numerical_deepbsde}
    为了与基于神经网络训练的 BSDE 方法比较,
    本文在 HJB-LQ 和 Allen--Cahn 两个半线性方程上记录了 DeepBSDE 的训练过程.
    DeepBSDE 的横轴采用实际运行时间; RFM 方法只有一次 LM 求解结果, 因而在图中以单点表示.

    \begin{table}
      \centering
      \caption{RFM 与 DeepBSDE 在半线性方程上的结果对比}
      \label{tab:deepbsde_compare}
      \begin{tabular}{llccccc}
        \toprule
        方程 & 方法 & 时间 / s & $y_0$ & 参考值 & 绝对误差 & 训练损失 \\
        \midrule
        HJB-LQ & RFM & 3.436 & 4.590000 & 4.590162 & $1.62\times10^{-4}$ & $9.80\times10^{-3}$ \\
        HJB-LQ & DeepBSDE & 15.221 & 4.586627 & 4.590162 & $3.53\times10^{-3}$ & 0.0260 \\
        Allen--Cahn & RFM & 3.145 & 0.053131 & 0.052785 & $3.46\times10^{-4}$ & $1.24\times10^{-5}$ \\
        Allen--Cahn & DeepBSDE & 29.192 & 0.052846 & 0.052785 & $6.14\times10^{-5}$ & 0.0020 \\
        \bottomrule
      \end{tabular}
    \end{table}

    在 HJB-LQ 方程上, RFM 在更短时间内得到更小的初值误差和训练损失.
    在 Allen--Cahn 方程上, DeepBSDE 经过更长训练后初值误差略小,
    但 RFM 的训练损失和运行时间仍明显更低.
    图~\ref{fig:deepbsde_rfm_comparison} 展示了训练损失和初值误差随时间变化的过程.

    \begin{figure}[ht]
      \centering
      \includegraphics[width=0.90\textwidth]{deepbsde-rfm-comparison}
      \caption{RFM 与 DeepBSDE 的训练损失和初值误差对比}
      \label{fig:deepbsde_rfm_comparison}
    \end{figure}

  \section{本章小结}\label{sec:numerical_summary}
    本章的四个算例覆盖了线性最小二乘和半线性 LM 两种求解情形.
    在线性方程中, Black--Scholes 与 Heat 方程的初值误差均约为 $10^{-4}$.
    在半线性方程中, HJB-LQ 和 Allen--Cahn 也能在标准配置下得到接近参考值的初值.
    样本数实验说明, 训练残差过小并不必然代表泛化良好, 因而独立测试残差是必要指标.
    维度实验显示, 当前实现在线性方程上的时间增长较慢, 半线性方程的时间随维度增长更快.
    与 DeepBSDE 的对比表明, 在本文两个半线性算例上, 随机特征方法能以较少训练时间得到有竞争力的初值精度.
